\section*{الفصل \ref{ch:g10:layered-ocean} --- المحيط الطبقي}

\begin{solution}{exo:g10:layered-ocean:1}
\omterm{def:g10:layered-ocean:mixed-layer}{الطبقة المختلطة} هي الطبقة السطحية من المحيط التي لا تكاد درجة الحرارة والملوحة
تتغيران فيها مع العمق؛ \omterm{def:g10:layered-ocean:thermocline}{وطبقة التدرج الحراري} هي الطبقة الواقعة تحتها التي تنخفض
فيها درجة الحرارة سريعًا مع العمق. وتقلّب الرياح والأمواج \omterm{def:g10:layered-ocean:mixed-layer}{الطبقة المختلطة} (ويساعدها
في الشتاء تبرد السطح).
\end{solution}

\begin{solution}{exo:g10:layered-ocean:2}
\omterm{def:g10:layered-ocean:upwelling}{الانبثاق} هو صعود ماء بارد من تحت الطبقة السطحية، ليحل محل ماء سطحي دفعته الريح
بعيدًا. ويجلب الماء الصاعد مغذيات معدنية إلى الطبقة المضاءة، حيث تستعملها العوالق
النباتية لتنمو؛ والعوالق النباتية تغذي بقية السلسلة الغذائية.
\end{solution}

\begin{solution}{exo:g10:layered-ocean:3}
المحيط الهادئ المداري: نحو \qty{29}{\degreeCelsius} عند السطح
و\qty{4.5}{\degreeCelsius} عند \qty{1000}{m}. والمحيط المتجمد الشمالي: نحو
\qty{-1.4}{\degreeCelsius} عند السطح وقرابة
\qty{0}{\degreeCelsius} عند \qty{1000}{m}.
\end{solution}

\begin{solution}{exo:g10:layered-ocean:4}
عند \qty{0}{\degreeCelsius} (\qty{1028.1}{kg/m^3} مقابل \qty{1023.3}{kg/m^3}
عند \qty{25}{\degreeCelsius}). وماء البحر: فعند \qty{4}{\degreeCelsius} تبلغ كثافة
الماء النقي \qty{1000.0}{kg/m^3}، وكثافة ماء البحر الذي ملوحته 35 نحو
\qty{1027.8}{kg/m^3} (مقروءة على المخطط).
\end{solution}

\begin{solution}{exo:g10:layered-ocean:5}
تنجرف العوامة على عمق \qty{1000}{m}، وتغوص إلى \qty{2000}{m}، ثم تصعد إلى
السطح وهي تقيس في طريقها درجة الحرارة والملوحة (من الناقلية) والضغط؛ وعند السطح
ترسل مقطعها عبر الأقمار الصناعية ثم تغوص من جديد. وتنجرف العوامات مع التيارات في
المحيط كله وتعمل سنوات، في كل الفصول، بعيدًا عن طرق السفن.
\end{solution}

\begin{solution}{exo:g10:layered-ocean:6}
$361.84\times10^6 \times 3.682 \approx \qty{1.33e9}{km^3}$، أي نحو
\num{1.33} مليار كيلومتر مكعب. $100/3682 \approx 0.027$: أي أقل من
\qty{3}{\percent} من متوسط العمق.
\end{solution}

\begin{solution}{exo:g10:layered-ocean:7}
$22.7 - 0.5 = \qty{22.2}{\degreeCelsius}$، وتُبلغ عند \qty{15}{m}: فعمق \omterm{def:g10:layered-ocean:mixed-layer}{الطبقة
المختلطة} نحو \qty{15}{m}. وبين \qty{20}{m} و\qty{50}{m}:
$(21.8 - 17.3)/30 = \qty{0.15}{\degreeCelsius}$ لكل متر.
\end{solution}

\begin{solution}{exo:g10:layered-ocean:8}
لا تسخّن الشمس إلا الطبقة السطحية، ولا ينتشر التسخين والتبريد الموسميان إلا عبر
\omterm{def:g10:layered-ocean:mixed-layer}{الطبقة المختلطة}: فالماء تحت \omterm{def:g10:layered-ocean:thermocline}{طبقة التدرج الحراري} معزول عنهما. وفي الشتاء يبرد السطح،
فيصير أكثف ويغوص في الماء الذي تحته، وتقلّب العواصف البحر: فيبلغ المزج عمقًا كبيرًا.
وفي الصيف تسخّن الشمس طبقة سطحية خفيفة تطفو على الماء الأبرد؛ ولا تقدر الرياح الألطف
على دفعها نحو الأسفل، فتبقى \omterm{def:g10:layered-ocean:mixed-layer}{الطبقة المختلطة} رقيقة.
\end{solution}

\begin{solution}{exo:g10:layered-ocean:9}
$0.78/0.18 \approx 4.3$ درجات قرب \qty{10}{\degreeCelsius}؛
و$0.78/0.06 = 13$ درجة قرب \qty{0}{\degreeCelsius}. ففي البحار القطبية الباردة
تؤثر التغيرات الصغيرة في الملوحة في الكثافة أكثر بكثير من تغيرات درجة الحرارة:
فالملوحة هناك هي التي تتحكم في \omterm{def:g10:layered-ocean:stratification}{التطبق}.
\end{solution}

\begin{solution}{exo:g10:layered-ocean:10}
درجة الحرارة واحدة: فالماء الأعمق أكثف لأنه أشد ملوحة (33.3 مقابل 28.1). والطبقة
الواقعة بينهما، حيث تتغير الملوحة سريعًا مع العمق، \omterm{def:g10:layered-ocean:thermocline}{طبقة تدرج ملحي}؛ وهي هنا أيضًا
\omterm{def:g10:layered-ocean:thermocline}{طبقة التدرج الكثافي}.
\end{solution}

\begin{solution}{exo:g10:layered-ocean:11}
حين تُدفع الكتلة نحو الأسفل تحتفظ بدرجة حرارتها وملوحتها، فبكثافتها. وإذ يحيط بها
ماء أكثف، تكون أخف من الماء الذي تزيحه: فتدفعها الطفوية عائدة إلى مستواها. والكتلة
المدفوعة نحو الأعلى تكون أكثف من محيطها الجديد فتغوص عائدة. فالطبقات مستقرة.
\end{solution}

\begin{solution}{exo:g10:layered-ocean:12}
الكتلة: $0.001 \times 1028.1 = \qty{1.0281}{kg}$. والحجم الجديد: $1.0281/1023.3
\approx \qty{0.0010047}{m^3}$، أي نحو \qty{1004.7}{mL}: فقد تمدد بنحو
\qty{4.7}{mL}، أي \qty{0.47}{\percent}. وطبقة سماكتها \qty{700}{m} كانت ستزداد
سماكتها بمقدار $700 \times 0.0047 \approx \qty{3.3}{m}$.
\end{solution}

\begin{solution}{exo:g10:layered-ocean:13}
في شمال الأطلسي في فبراير. فهناك تكون لأعلى \qty{150}{m} درجة حرارة شبه واحدة،
فكثافة شبه واحدة: وتقليبها لا يكلف إلا قليلًا. أما في المحيط الهادئ المداري فترتفع
الكثافة بمقدار $1026.0 - 1021.6 =
\qty{4.4}{kg/m^3}$ بين \qty{50}{m} و\qty{200}{m}: وهي طبقة تدرج كثافي قوية،
وكان على العاصفة أن تدفع ماء خفيفًا نحو الأسفل داخل ماء أكثف منه بكثير.
\end{solution}

\begin{solution}{exo:g10:layered-ocean:14}
يخضور مرتفع بعيدًا عن السواحل: الشريط الممتد على خط الاستواء في المحيط الهادئ
(وفي الأطلسي)، حيث تدفع الرياح التجارية انبثاقًا؛ والعروض العليا، كشمال الأطلسي،
حيث تبلغ \omterm{def:g10:layered-ocean:mixed-layer}{الطبقة المختلطة} الشتوية عمقًا كبيرًا فترفع المغذيات. ويخضور شديد الانخفاض:
مراكز الدوامات شبه المدارية، الشديدة \omterm{def:g10:layered-ocean:stratification}{التطبق}، حيث لا شيء يرفع الماء العميق وقد
نفدت المغذيات من الطبقة المضاءة.
\end{solution}

\begin{solution}{exo:g10:layered-ocean:15}
الغوص: $1000/(6 \times 3600) \approx \qty{0.046}{m/s}$، أي نحو
\qty{4.6}{cm/s}. والصعود: $2000/(10 \times 3600) \approx \qty{0.056}{m/s}$،
أي نحو \qty{5.6}{cm/s}. $36\,000/2 = 18\,000$ قياس، يفصل بينها في المتوسط
$2000/18\,000 \approx \qty{0.11}{m}$.
\end{solution}

\begin{solution}{pb:g10:layered-ocean:1}
\textbf{1.} الدوامة، بفارق $24.0 - 19.0 = \qty{5.0}{\degreeCelsius}$.

\textbf{2.} الدوامة (\qty{20.5}{\degree} ج مقابل \qty{13.5}{\degree} ج).
فهي أبعد عن خط الاستواء، فكان ينبغي أن تكون الأبرد: ولا يمكن لخطي العرض وحدهما
أن يفسّرا لماذا هي أدفأ بمقدار \qty{5}{\degreeCelsius}.

\textbf{3.} $24.0 - 0.5 = \qty{23.5}{\degreeCelsius}$، وتُبلغ بين
\qty{50}{m} (23.6) و\qty{100}{m} (22.6)، عند نحو $50 + 50 \times
0.1/1.0 = \qty{55}{m}$.

\textbf{4.} لا: فدرجة الحرارة تنخفض باطراد من السطح، بمقدار
\qty{1.6}{\degreeCelsius} في أول \qty{25}{m} و\qty{5.3}{\degreeCelsius}
في أول \qty{100}{m}؛ \omterm{def:g10:layered-ocean:thermocline}{فطبقة التدرج الحراري} تبدأ من السطح.

\textbf{5.} بين \qty{200}{m} و\qty{300}{m}: $20.3 - 14.4 =
\qty{5.9}{\degreeCelsius}$ لكل مئة متر.

\textbf{6.} قبالة بيرو، بين \qty{50}{m} و\qty{100}{m}: نحو
\qty{70}{m} بالاستيفاء (ويعطي المقطع الكامل \qty{65}{m}). وفي
الدوامة، بين \qty{200}{m} و\qty{300}{m}: نحو $200 + 100 \times 5.3/5.9
\approx \qty{290}{m}$.

\textbf{7.} نعم، عند كل مستوى في العمودين: فالماء الأخف يستقر فوق الماء الأكثف،
والعمودان كلاهما مستقران.

\textbf{8.} كان ينبغي للماء الأشد ملوحة أن يكون أكثف؛ ومع ذلك فالماء السطحي في الدوامة
أخف: \qty{1024.71}{kg/m^3} مقابل \qty{1025.08}{kg/m^3}.

\textbf{9.} ماء الدوامة أشد ملوحة بنحو 1.3 وحدة، وهذا يضيف نحو
$1.3 \times 0.8 \approx \qty{1.0}{kg/m^3}$؛ وهو أدفأ بمقدار \qty{5}{\degreeCelsius}،
وهذا ينقص نحو $5 \times 0.27 \approx \qty{1.4}{kg/m^3}$ عند درجات الحرارة هذه.
فدرجة الحرارة هي الغالبة.

\textbf{10.} $1026.44 - 1025.08 = \qty{1.36}{kg/m^3}$.

\textbf{11.} $1025.30 - 1024.71 = \qty{0.59}{kg/m^3}$.

\textbf{12.} من السطح إلى \qty{500}{m} ترتفع الكثافة بمقدار
$1026.92 - 1024.71 = \qty{2.21}{kg/m^3}$، لا يقع منها فوق \qty{200}{m} إلا
\qty{0.59}{kg/m^3}: فمعظمها، \qty{1.62}{kg/m^3}، يقع تحت
\qty{200}{m}، في \omterm{def:g10:layered-ocean:thermocline}{طبقة التدرج الحراري} بين \qty{200}{m} و\qty{500}{m}.

\textbf{13.} قبالة بيرو ترتفع الكثافة بمقدار $1026.97 - 1025.08 =
\qty{1.89}{kg/m^3}$ من السطح إلى \qty{500}{m}، منها
\qty{1.36}{kg/m^3} في أعلى \qty{200}{m}: فمعظمها قرب السطح.

\textbf{14.} عند \qty{500}{m} يتشابه الموقعان تقريبًا
(\qty{7.9}{\degreeCelsius} و\qty{7.1}{\degreeCelsius}،
\qty{1026.97}{kg/m^3} و\qty{1026.92}{kg/m^3}): فالماء البارد نفسه يقع
تحتهما كليهما؛ والفرق في مدى اقترابه من السطح.

\textbf{15.} الريح الهابة على امتداد الشاطئ نحو خط الاستواء تدفع الماء السطحي نحو
عرض البحر؛ فيصعد ماء من الأسفل قرب الساحل ليحل محله. فيُسحب الماء البارد \omterm{def:g10:layered-ocean:thermocline}{لطبقة
التدرج الحراري} نحو الأعلى، فتميل \omterm{def:g10:layered-ocean:thermocline}{طبقة التدرج الحراري} صاعدة نحو الساحل حتى تكاد
تبلغ السطح.

\textbf{16.} قبالة بيرو: الماء عند \qty{13}{\degreeCelsius} إلى
\qty{15}{\degreeCelsius}، القادم من تحت \omterm{def:g10:layered-ocean:thermocline}{طبقة التدرج الحراري}، يقع ضمن أعلى
\qty{100}{m}، في متناول تقليب الريح والضوء. أما في
الدوامة فيقع تحت \qty{250}{m}، تحت طبقة دافئة خفيفة لا تقدر الريح على تقليبها إلى
هذا العمق.

\textbf{17.} كان \omterm{def:g10:layered-ocean:upwelling}{الانبثاق} سيتوقف: فيسخن الماء السطحي في الشمس، وتتكون طبقة دافئة
خفيفة في الأعلى، ويتضاءل إمداد الطبقة المضاءة بالمغذيات.

\textbf{18.} ساحل بيرو أزرق فاتح إلى أصفر (يخضور كثير)؛ ووسط دوامة جنوب المحيط
الهادئ من أشد مناطق الخريطة زرقة (قليل جدًا).

\textbf{19.} تحتاج العوالق النباتية إلى المغذيات كما تحتاج إلى الضوء. وفي الدوامة
يُبقي \omterm{def:g10:layered-ocean:stratification}{التطبق} القوي الماء العميق الغني بالمغذيات بعيدًا تحت الطبقة
المضاءة، وقد استُهلكت مغذيات الطبقة السطحية منذ زمن بعيد.

\textbf{20.} $1.36/0.59 \approx 2.3$: \textbf{ازدياد الكثافة في أعلى
\qty{200}{m} أكبر قبالة بيرو منه في الدوامة بنحو 2.3 مرة}. \omterm{def:g10:layered-ocean:thermocline}{فطبقة
التدرج الكثافي}، والماء البارد الغني بالمغذيات تحتها، رُفعت إلى أعلى \qty{200}{m}
قبالة بيرو، بينما تقع في الدوامة أعمق، تحت
\qty{200}{m}.
\end{solution}
